%! The Normal Distribution and z-Scores — Unit 1, Lesson 1.7 — Experience & Formalize
%! (Activity · QuickNotes · Application) (Answer Key)
% EFFL component, Math Medic sizing: 12pt body, OPEN white space after every question.
% Page budget: Activity <=2pp · QuickNotes 1/2pp · Application 1/2-1pp — THREE pages total.
% Check Your Understanding is NOT in this component: those items were moved to the end of
% the packet and are now the lesson's SCORED homework (see ../homework/).
%
% WRITE-SPACE BUDGET IS FIXED — DO NOT SHAVE IT. Activity = 11.2cm across 7 sub-questions:
%   1a 1.4 · 1b 1.4 · 1c 2.2 · 1d 1.2 · 2b 1.2 · 2c 2.2 · 2d 1.6
% Application = 5.0cm: (a) 1.4 · (b) 1.4 · (c) 2.2.  1.2-2.2cm is the usable floor at 12pt.
%
% HOW PROBLEM 2 IS BUILT, AND WHY (this shape is what got the file from 4pp to 3pp; prose
% compression was exhausted before it):
%   * ONE merged results table, rows = the three bands, columns = BOTH contexts side by side.
%     This is pedagogy first and compression second: "the same three numbers, twice" is only
%     visible if 68.2 sits next to 68.25 on the page instead of across a page turn.
%   * ONE band-count strip, three rows (band label / loaves / readings). The per-context weight
%     and temperature RANGE rows were removed — the ranges survive where they are actually
%     used: on the histogram's axis for the bread, and in 2d's "more than 840 g".
%   * ONE histogram (the bread). The clinic is table-only ON PURPOSE — students infer its shape
%     from its counts, which reinforces that the empirical rule is about the shape of the data
%     and not about having a picture. The bread histogram keeps its band-boundary ticks and its
%     dashed mean line; those are doing work.
%   * 2b's predict-before-you-look step survives the merge: the table shows the clinic's raw
%     COUNTS, never its percentages, so the prediction cannot be answered by reading across —
%     and the stem says "do not compute yet" explicitly.
% \boxguard values on the last three boxes are 12/12/10 rather than 14: at 14 the Application
% would not start at the foot of p.3 and was thrown whole onto p.4.
%
% SPOILER RULE: the Activity never says parameter, mu, sigma, normal, normal curve,
% approximately normal, empirical rule, 68-95-99.7, standardized score, or z-score. It says
% "the average of the field", "the spread", "how many spreads away". Mean, standard deviation,
% percentile, IQR and the five-number summary WERE formalized in Lessons 1.4-1.6 and are prior
% knowledge here, so `standard deviation' is used freely.
%
% NO SKETCH QUESTIONS (hard constraint): every dotplot and histogram is pre-drawn. Students
% "construct" only by filling in counts and percents, never by drawing an axis or a curve.
%
% THE NUMBERS
%   Problem 1 — shot put field (20 throws, m): 6.0 7.0 7.0 7.5 8.0 8.0 8.0 8.5 8.5 9.0 9.0 9.0
%     9.5 9.5 10.0 10.0 10.5 11.0 11.5 12.0  -> mean 8.975 ~ 9.0, s ~ 1.53 ~ 1.5. Maya = 12.0.
%   100 m field (20 times, s): 12.3 12.5 12.6 12.7 12.8 12.8 12.9 13.0 13.1 13.2 13.2 13.3
%     13.4 13.5 13.6 13.6 13.7 13.9 14.1 14.5 -> mean 13.235 ~ 13.2, s ~ 0.56 ~ 0.6. Dev = 12.3.
%     Maya: (12.0-9.0)/1.5 = +2.0.   Dev: (12.3-13.2)/0.6 = -1.5, and LOWER IS BETTER.
%   Problem 2 — bread, 500 loaves, mean 800 g, spread 20 g. Bands 11/68/170/171/67/12 (+1
%     outside): within 1 = 341 (68.2%), within 2 = 476 (95.2%), within 3 = 499 (99.8%).
%   Temperature, 400 readings, mean 98.2 F, spread 0.7 F. Bands 9/54/137/136/54/9 (+1 outside):
%     within 1 = 273 (68.25%), within 2 = 381 (95.25%), within 3 = 399 (99.75%).
%   The two sets of percentages agreeing is the second discovery (EK VAR-2.A.3).
%
% Every \begin{work} block is authored BYTE-IDENTICALLY here and in experience/ — under
% -boxes each ships as a \vphantom of exact height, under -key it prints. They cannot drift.
% GENERATED FROM experience/main.tex — the blank with answers filled in, so every reserved
% height, guard, scale, \vspace, \arraystretch and itemsep is identical to it. NO teachernote
% here — that lives in the lesson plan.
\documentclass[12pt]{article}
\usepackage{apstats-article}
\usepackage{apstats-key}   % pulls in -boxes; do NOT also load -boxes
\newcommand{\answerspace}[2]{\par\nopagebreak\noindent\begin{minipage}[t][#1][t]{\linewidth}%
  \color{keyred}\bfseries #2\end{minipage}\par}

\begin{document}

\pageheader{Unit 1, Lesson 1.7}{Experience \& Formalize: Two Tryouts --- Answer Key}
% NAMESTRIP: no \namedateperiod here — mirrors the blank.

\vspace{0.02in}

% ── Part 1: the Activity ─────────────────────────────────────────────────────
\begin{headlinebox}{sky}
\textbf{Two Tryouts.} One scholarship, two candidates, two different events --- and somebody has
to defend the decision to the athlete who loses.
\end{headlinebox}

\vspace{0.03in}

\begin{scenariobox}[1. Who Had the Better Day?]{navy}
\small
Maya throws the shot put; Dev runs the 100 m dash. Every athlete in each event at Saturday's
meet, with a summary of each field:

\vspace{3pt}
\begin{center}
\begin{tikzpicture}[scale=0.70]
  % ── left: shot put field, one dot per throw ──
  \draw[->, thick] (-0.25,0) -- (6.3,0);
  \foreach \v in {6,7,8,9,10,11,12} {
    \draw ({(\v-5.5)*0.85},-0.08) -- ({(\v-5.5)*0.85},0.08);
    \node[below, font=\scriptsize] at ({(\v-5.5)*0.85},-0.12) {\v};}
  \foreach \v/\n in {6.0/1, 7.0/2, 7.5/1, 8.0/3, 8.5/2, 9.0/3, 9.5/2, 10.0/2, 10.5/1, 11.0/1,
    11.5/1} {
    \foreach \k in {1,...,\n} { \fill[navy] ({(\v-5.5)*0.85},{0.2*\k}) circle (0.075); }}
  \fill[goldacc] ({(12.0-5.5)*0.85},0.2) circle (0.1);
  \node[font=\scriptsize\bfseries, text=goldacc] at ({(12.0-5.5)*0.85},0.6) {Maya};
  \node[font=\scriptsize] at (2.9,-0.55) {Distance thrown (metres)};
  \node[font=\bfseries\small] at (2.9,1.15) {Shot put: 20 throwers};
  % ── right: 100 m field, one dot per runner ──
  \begin{scope}[xshift=8.0cm]
    \draw[->, thick] (-0.25,0) -- (6.2,0);
    \foreach \v/\lab in {12.5/12.5, 13.0/13.0, 13.5/13.5, 14.0/14.0, 14.5/14.5} {
      \draw ({(\v-12.2)*2.4},-0.08) -- ({(\v-12.2)*2.4},0.08);
      \node[below, font=\scriptsize] at ({(\v-12.2)*2.4},-0.12) {\lab};}
    \foreach \v/\n in {12.5/1, 12.6/1, 12.7/1, 12.8/2, 12.9/1, 13.0/1, 13.1/1, 13.2/2, 13.3/1,
      13.4/1, 13.5/1, 13.6/2, 13.7/1, 13.9/1, 14.1/1, 14.5/1} {
      \foreach \k in {1,...,\n} { \fill[navy] ({(\v-12.2)*2.4},{0.2*\k}) circle (0.075); }}
    \fill[goldacc] ({(12.3-12.2)*2.4},0.2) circle (0.1);
    \node[font=\scriptsize\bfseries, text=goldacc] at ({(12.3-12.2)*2.4},0.6) {Dev};
    \node[font=\scriptsize] at (2.9,-0.55) {Finishing time (seconds)};
    \node[font=\bfseries\small] at (2.9,1.15) {100 m dash: 20 runners};
  \end{scope}
\end{tikzpicture}
\end{center}
\vspace{2pt}

\renewcommand{\arraystretch}{1.1}
\begin{tabularx}{\linewidth}{@{} l Y Y Y @{}}
\rowcolor{sky}
\textbf{Event} & \textbf{Field average} & \textbf{Field standard deviation} &
\textbf{Our athlete} \\
\hline
Shot put (metres) & 9.0 m & 1.5 m & Maya: 12.0 m \\
100 m dash (seconds) & 13.2 s & 0.6 s & Dev: 12.3 s \\
\end{tabularx}

\vspace{2pt}
\normalsize
\begin{enumerate}[label=\textbf{\alph*.}, leftmargin=1.8em, itemsep=4pt]
  \item Maya's result is 12.0, Dev's is 12.3. Whose \emph{number} is bigger? \ans{Dev's (12.3)}
        Does that settle it? Say why not.
        \answerspace{1.4cm}{No. One is metres thrown, the other seconds run --- different scales,
        and not even the same kind of measurement.}
  \item How far is each result from that athlete's \emph{own} event average?
        \begin{work}
        \text{Maya} &= 12.0 - 9.0 = +3.0 \text{ m} \qquad
        \text{Dev} = 12.3 - 13.2 = -0.9 \text{ s}
        \end{work}
        Does \emph{that} settle it? \ans{No} Say why it still does not.
        \answerspace{1.4cm}{The differences are still in different units: 3.0 metres against
        0.9 seconds. Nothing says which is the bigger achievement.}
  \item \textbf{Now use each event's own standard deviation from the table.} How many of their
        own event's standard deviations does each result sit from that event's average?
        \begin{work}
        \text{Maya} &= \frac{3.0 \text{ m}}{1.5 \text{ m}} = 2.0 \qquad
        \text{Dev} = \frac{-0.9 \text{ s}}{0.6 \text{ s}} = -1.5
        \end{work}
        Whose result is more unusual \emph{for their own event}? In a sprint, is \emph{below}
        average good or bad --- does that change how you read Dev's number?
        \answerspace{2.2cm}{Maya --- 2.0 spreads from her event's average against Dev's 1.5. The
        units cancel, so both are plain counts and can be compared. Below average is \emph{good}
        in a sprint: Dev beat his field, just not by as unusual a margin.}
  \item Jordan cleared exactly the average height for the high-jump field. What does part (c)'s
        method give for Jordan? \ans{0} What would a \emph{negative} result mean?
        \answerspace{1.2cm}{That the result is below that event's average, by that many of the
        event's own standard deviations.}
\end{enumerate}
\end{scenariobox}

\boxguard[12]
\begin{scenariobox}[2. The Same Three Numbers, Twice]{navy}
\small
Two unrelated records: a bakery's \textbf{500} loaves and a clinic's \textbf{400} body
temperatures, each sorted into bands one standard deviation wide from its own average. (One loaf
and one reading fell outside.) The loaves are pictured; the clinic's counts are in the table.

\vspace{3pt}
\begin{center}
\begin{tikzpicture}[scale=0.70]
  \foreach \i/\c in {0/3, 1/8, 2/24, 3/44, 4/74, 5/96, 6/96, 7/75, 8/43, 9/24, 10/8, 11/4} {
    \fill[navy] ({\i*0.5},0) rectangle ({(\i+1)*0.5},{\c*0.019});
    \draw[white, very thin] ({\i*0.5},0) rectangle ({(\i+1)*0.5},{\c*0.019});}
  \draw[->, thick] (0,0) -- (0,2.3);
  \draw[->, thick] (0,0) -- (6.4,0);
  \foreach \y/\lab in {0/0, 0.95/50, 1.90/100} {
    \draw (-0.08,\y) -- (0.08,\y); \node[left, font=\tiny] at (-0.1,\y) {\lab};}
  \foreach \x/\lab in {0/740, 1/760, 2/780, 3/800, 4/820, 5/840, 6/860} {
    \draw (\x,-0.08) -- (\x,0.08); \node[below, font=\tiny] at (\x,-0.1) {\lab};}
  \draw[redacc, dashed, thin] (3,0) -- (3,2.0);
  \node[font=\tiny, text=redacc] at (3,2.18) {average 800 g};
  \node[font=\scriptsize] at (3,-0.55) {Weight of a loaf (grams)};
\end{tikzpicture}
\end{center}
\vspace{2pt}

\renewcommand{\arraystretch}{1.05}
\begin{tabularx}{\linewidth}{@{} l Y Y Y Y Y Y @{}}
\rowcolor{sky}
\textbf{Spreads from average} & $-3$ to $-2$ & $-2$ to $-1$ & $-1$ to 0 & 0 to $+1$ &
  $+1$ to $+2$ & $+2$ to $+3$ \\
\hline
\textbf{Loaves} (500 total) & 11 & 68 & 170 & 171 & 67 & 12 \\
\textbf{Readings} (400 total) & 9 & 54 & 137 & 136 & 54 & 9 \\
\end{tabularx}

\vspace{4pt}
\begin{tabularx}{\linewidth}{@{} l Y Y Y Y @{}}
\rowcolor{sky}
\textbf{Within \ldots{} of the average} & \textbf{Loaves} & \textbf{Percent of 500} &
\textbf{Readings} & \textbf{Percent of 400} \\
\hline
1 standard deviation & \ans{341} & \ans{68.2\%} & \ans{273} & \ans{68.25\%} \\
2 standard deviations & \ans{476} & \ans{95.2\%} & \ans{381} & \ans{95.25\%} \\
3 standard deviations & \ans{499} & \ans{99.8\%} & \ans{399} & \ans{99.75\%} \\
\end{tabularx}

\vspace{2pt}
\normalsize
\begin{enumerate}[label=\textbf{\alph*.}, leftmargin=1.8em, itemsep=4pt]
  \item \textbf{The bread.} Add the counts from the strip above and fill in the \emph{Loaves} and
        \emph{Percent of 500} columns.
  \item \textbf{Predict first --- do not compute yet.} Cover the clinic's row. What do you think
        its three percentages will be? Write them here, then fill in the last two columns.
        \answerspace{1.2cm}{Predictions vary. Accept anything written down before the clinic's
        columns were filled in.}
  \item \textbf{The two sets of percentages agree to within a fraction of a percent}, in two
        records with nothing in common. What do the two distributions share that explains it ---
        and what would make a data set \emph{not} come out this way?
        \answerspace{2.2cm}{Both are mound-shaped and symmetric about one centre --- the clinic's
        counts rise and fall just as the loaves' picture does. Shape, not units or spread, is
        what fixes the percentages. A skewed, two-peaked, or flat set would not.}
  \item Using \emph{only} your part (a) percentages: roughly what fraction of the loaves weigh
        more than 840 g, two spreads above average? Explain the arithmetic.
        \answerspace{1.6cm}{About 2.5\%. Roughly 95\% are within two spreads, so about 5\% are
        outside; symmetry splits that evenly, leaving 2.5\% above 840 g.}
\end{enumerate}
\end{scenariobox}

% ── Part 2: QuickNotes (the debrief fills this) — target: half a page ────────
\boxguard[12]
\begin{tcolorbox}[enhanced, colback=sky, colframe=navy, boxrule=0.8pt, arc=2mm,
    left=3mm, right=3mm, top=1.5mm, bottom=1.5mm, breakable,
    fonttitle=\bfseries\small\color{white}, title={QuickNotes: The Normal Distribution and z-Scores},
    attach boxed title to top left={yshift=-2mm, xshift=4mm},
    boxed title style={colback=navy, colframe=navy, arc=1.5mm}]
\small
\begin{itemize}[leftmargin=1.1em, itemsep=2pt]
  \item \textbf{Parameter:} a numerical summary of a \ans{population}. \textbf{Statistic:} a
        numerical summary of a \ans{sample}. The parameters \ans{$\mu$} and \ans{$\sigma$}
        describe a population; $\bar x$ and $s_x$ (Lesson 1.5) describe a sample.
  \item \textbf{Normal distribution:} a curve that is \ans{mound-shaped} and \ans{symmetric}, fixed
        entirely by $\mu$ and $\sigma$. Real data is never exact, so we say \ans{approximately} normal.
  \item \textbf{Empirical rule.} Approximately normal $\Rightarrow$ about \ans{68\%} of
        observations within 1 standard deviation of the mean, \ans{95\%} within 2,
        \ans{99.7\%} within 3.
  \item \textbf{Standardized score ($z$-score).} $z = (x_i - \mu)/\sigma$ counts how many
        \ans{standard deviations} the value sits from the mean. In context: ``a 830 g loaf has $z = 1.5$, so
        it is \ans{1.5 standard deviations} above the mean weight of 800 g.''
  \item $z < 0$: the value is \ans{below} the mean. $z = 0$: \ans{exactly at} it. Carrying no
        units, $z$-scores and \ans{percentiles} report \emph{relative position}, so they compare two
        different data sets.
  \item Technology (calculator or normal table) turns a $z$-score into an exact \ans{proportion}.
        That is Unit~5.
\end{itemize}
\end{tcolorbox}

% ── Part 3: the Application (worked together, ~6 min) ────────────────────────
\boxguard[10]
\begin{notesbox}{Application: The Bread Loaf}
Together. Back to the bakery: loaf weights are approximately normal, $\mu = 800$ g,
$\sigma = 20$ g.

\begin{enumerate}[label=\textbf{\alph*.}, leftmargin=1.8em, itemsep=3pt]
  \item A loaf weighs \textbf{834 g}. Find its $z$-score, then interpret it in one sentence, in
        context.
        \begin{work}
        z &= \frac{x_i - \mu}{\sigma} = \frac{834 - 800}{20} = \frac{34}{20} = 1.7
        \end{work}
        \answerspace{1.2cm}{This loaf weighs 1.7 standard deviations more than the mean loaf
        weight of 800 grams --- heavier than typical, but not extreme.}
  \item By the empirical rule, about what percent of loaves weigh \emph{more than} 840 g?
        \ans{About 2.5\%} Explain how.
        \answerspace{1.2cm}{840 g is two standard deviations above 800 g; about 5\% of loaves lie
        outside two, and symmetry puts half of that above 840 g.}
  \item A customer returns a \textbf{754 g} loaf, calling it ``way underweight.''
        \begin{work}
        z &= \frac{754 - 800}{20} = \frac{-46}{20} = -2.3
        \end{work}
        Is that fair? Judge it in context, and name one limit of the model you used.
        \answerspace{2.0cm}{Broadly fair: at $z = -2.3$ the loaf sits more than two standard
        deviations below the 800 g mean, so under about 2.5\% of loaves are that light. But the
        normal model is only approximate, and one loaf says nothing about the rest.}
\end{enumerate}
\end{notesbox}

\end{document}
